\section{模型评价与推广}

\subsection{模型优点}

\begin{enumerate}[label=(\arabic*), leftmargin=2.2em, itemsep=2pt]
\item \textbf{机理完备、物理可行}。模型显式刻画了逐时段功率平衡（并允许弃光，避免高光伏低负荷日无解）、
储能 SOC 递推（单程 90\% 效率、往返 81\%）、充放电功率与 SOC 双向限幅。
全部解经物理复检：功率平衡残差 $\le1.1\times10^{-13}$ kWh，SOC 始终位于 $[1200,10800]$，
无同时充放电时段，充放电均未越限。
\item \textbf{决策全部由 LP 精确求解}。四问共约 4000 个线性规划（全年 365 天 $\times$ 多场景），
均在秒级内求得全局最优；对照实验证明启发式在本问题中明显劣于 LP（典型日贵 20.1\%）。
\item \textbf{预测模型是本题最大的杠杆，且用“可分解 + 可堆叠”的方式实现}：
负载用“水平$\times$形状分解 + 周五六/其他分类凸组合”（MAE 137.5 $\to$ 132.3 kW，
其中周五六 $-8.0\%$、其余五天 $-1.8\%$），光伏用“因果岭回归堆叠”
（MAE 152.0 $\to$ 136.9 kW，$-9.9\%$），电价用“廓线+自适应水平+GBDT 堆叠”
（MAE 0.0581 $\to$ 0.0367 元/kWh）；全年费用由 13\,848\,183 元降到 \textbf{13\,655\,835 元}。
\item \textbf{用预言机分解量化信息价值}。把距完全信息下界的 136.9 万元缺口分解为
光伏噪声 49.1 万元（35.9\%）、负荷噪声 42.4 万元（31.0\%）、交互 33.2\%，
据此判断“把预算投在实时平衡控制上比投在更复杂的预测模型上更划算”：
实时平衡控制相对僵硬执行省下 90.2 万元，与“完美预知全部噪声”的理论上限 91.5 万元同量级，
而后者不可能实现。
\item \textbf{计划与结算分开建模，保证策略可实行}。0:00 计划只用历史信息，
日内结算用第 \ref{sec:q2} 节的\textbf{实时平衡控制规则}（只用当前实测与当前储电量），
两层模型独立求解，决策者的信息集随时间单调增长，
从机制上排除了“用未来数据事后调整计划”的虚优（\cref{tab:q2exec} 给出同一计划三种执行方式的对照）。
\item \textbf{结果可复现、可核查}。全部结果由 40 余个 Python 脚本（提交材料“代码”部分）
从原始附件自动生成，输出的 5 个结果文件与附件5 模板逐单元格对齐；
论文中每个数字都可回溯到脚本与中间数据。另有两条独立的正确性证据：
①\textbf{审计程序}（\texttt{audit\_results.py}）不依赖任何中间变量，直接重读 5 个结果文件，
用附件2/附件4 的原始实测逐时段复算功率平衡、SOC 递推、充放电功率与 SOC 上下限、
四段计价费用，四问全部通过；
②\textbf{求解器稳健性对照}（\texttt{solver\_robustness\_final.py}）把同一条管道分别交给
HiGHS 对偶单纯形与内点法求解，问题二/问题三费用相对差 $<0.1\%$，
说明结论不依赖“返回了哪一个最优顶点”。
\end{enumerate}

\subsection{模型缺点与局限}

\begin{enumerate}[label=(\arabic*), leftmargin=2.2em, itemsep=2pt]
\item \textbf{滚动调整只用了“预报”，未利用“已观测”信息}。6:00 调整时，0:00--6:00 的实际
光伏与负载已经发生，若把这些观测纳入（如按“已观测 + 预报”分段重估），应可进一步降低紧急购电；
本文保守地只使用预报，使结论更稳健（不依赖额外假设）。
\item \textbf{紧急购电未设功率上限}。题面未给出外网联络线容量，本文按不限功率处理；
若存在容量约束，需在结算模型中加 $E^{\text{急}}_t\le \bar E$，此时高光伏波动日的
缺口可能无法完全消除，模型需引入切负荷变量并重新定义可行域。
\item \textbf{日末残值用“次日平均电价”近似}。严格的滚动时域应求解跨日耦合的多日联合优化；
本文用残值系数 $v=0.9\bar p$ 近似值函数，灵敏度分析表明该近似对结果影响 $<0.3\text{\%}$。
\item \textbf{储能参数未做容量-功率联合优化}。本文以附录1 给定的 12000 kWh / 5000 kW 为常量，
而实际规划中储能容量本身是可决策变量（投资成本与运行收益的权衡）。
\item \textbf{未考虑电池寿命与爬坡约束}。频繁浅充放会加速容量衰减，本文按理想储能处理。
\end{enumerate}

\subsection{模型改进与推广}

\begin{enumerate}[label=(\arabic*), leftmargin=2.2em, itemsep=2pt]
\item \textbf{把“已观测”信息纳入滚动决策}：在式 \eqref{eq:adj_obj}--\eqref{eq:adj_link} 中，
对已经过去的时段用实测值、对未发生的时段用预报值，可望进一步降低紧急购电；
当信息价值足够大时（可用本文的预言机分解先核算）再引入机会约束或分布鲁棒优化，
在“成本”与“鲁棒性”之间取得平衡。
\item \textbf{储能容量-功率联合优化}：把 $\bar C$、容量与 SOC 限幅作为投资决策变量，
目标改为“年运行费 + 年化投资费”，得到最优储能配置。
\item \textbf{需求侧响应与多微网协同}：把可平移负载、电动汽车、共享储能纳入模型，
目标函数扩展为多微网总成本最小（可用纳什谈判/分布式优化求解）。
\item \textbf{竞价为不确定量的场景}：若不允许假设“已知实时电价”，可用本文验证过的
“同类星期廓线 + 自适应水平 + GBDT 堆叠”电价预测模型；实测表明该假设的代价不但为 0，
用预测电价定计划反而使全年费用更省 0.20\%（第 \ref{sec:q4} 节）。
\item \textbf{方法学推广}：本文“\textbf{规律挖掘 $\to$ 依据构造 $\to$ LP 精确求解 $\to$
储能实时平衡控制 $\to$ 信息价值核算}”的范式，可直接迁移到园区综合能源系统、
数据中心“算电协同”、光储充一体化电站等同类问题。
\end{enumerate}
